\documentclass[11pt,a4paper]{article}
\usepackage[utf8]{inputenc}
\usepackage[T1]{fontenc}
\usepackage{amsmath,amssymb,amsfonts}
\usepackage{hyperref}
\usepackage{geometry}
\usepackage{enumitem}
\usepackage{booktabs}
\usepackage{xcolor}
\geometry{margin=1in}

\title{Briefing: Richard Goodman / apoth3osis — Meet-and-Greet}
\author{QWAV Strategy — Prepared May 2026}
\date{}

\begin{document}

\maketitle

\vspace{-2em}
\begin{center}
\small
\textbf{Date:} May 2026 \quad|\quad
\textbf{Status:} Pre-meeting briefing \quad|\quad
\textbf{Contact:} Richard Goodman, Apoth3osis Labs — \texttt{rgoodman@apoth3osis.io}
\end{center}

\hrulefill

\section{Briefing: Richard Goodman / apoth3osis --- Meet-and-Greet}


\textbackslash{}textbf\{Date:\} May 2026
\textbackslash{}textbf\{Purpose:\} Initial video call to explore formal verification collaboration for ultrametric/Bruhat-Tits quantum computing claims.
\textbackslash{}textbf\{Status:\} Richard initiated contact via QWAV website inquiry (April 2026). Email thread covers \textasciitilde{}6 weeks of substantive technical exchange.
\textbackslash{}textbf\{Primary Contact:\} Richard Goodman, Apoth3osis Labs --- \textbackslash{}texttt\{rgoodman@apoth3osis.io\}
\textbackslash{}textbf\{Website:\} https://www.apoth3osis.io --- "Paper $\to$ Proof $\to$ Code"

\hrulefill


\subsection{1. Meeting Purpose}


Explore whether Richard Goodman can formally verify core QWAV claims in Lean 4, specifically:

\begin{enumerate}
\item \textbackslash{}textbf\{The strong triangle inequality produces passive error confinement\} in Bruhat-Tits tree encodings
\item \textbackslash{}textbf\{A threshold theorem for ultrametric quantum error correction\}  --- analogous to the surface code threshold theorem from the 1990s, but geometric rather than measurement-based
\item \textbackslash{}textbf\{Energy barrier scaling:\} $E_{\text{barrier}} \propto q^d$ with tree depth
\end{enumerate}


\textbackslash{}textbf\{The strategic goal:\} Use formal verification (Lean 4 proofs) to bypass both peer review and physical lab experiments. A Lean proof that type-checks is unassailable  --- no reviewer, no gatekeeper, no credential check. This is the ultimate realization of the substance-first strategy.

\textbackslash{}textbf\{The practical goal:\} A co-authored paper with machine-verifiable repository proving specific QWAV claims as theorems, using known physical parameters from published quantum computing experiments (Google, IBM, Rigetti, etc.) rather than requiring new lab data.

\hrulefill


\subsection{2. Who Is Richard Goodman}


\begin{tabular}{|l|l|}
\hline
\textbf{Item} & \textbf{Detail} \\
\hline
\textbackslash{}textbf\{Role\} & Founder, Apoth3osis Labs \\
\hline
\textbackslash{}textbf\{Platform\} & HeytingLean  --- formal verification stack in Lean 4 \\
\hline
\textbackslash{}textbf\{Scale\} & 4,649 Lean 4 files, 876,000+ lines, \textbackslash{}textbf\{zero \textbackslash{}texttt\{sorry\}\} in core \\
\hline
\textbackslash{}textbf\{Pipeline\} & 25-phase verified compilation pipeline reducing Lean 4's type checker to SKY combinators (S, K, Y) \\
\hline
\textbackslash{}textbf\{SKY proof\} & 18,000 lines, zero non-standard axioms \\
\hline
\textbackslash{}textbf\{Architecture\} & Six-lens categorical transport: Laws of Form, Heyting, Tensor, Clifford, Graph, Geometric \\
\hline
\textbackslash{}textbf\{Collaborator\} & Vladimir Veselov (Russian mathematician, student of Kolmogorov)  --- ultrametric structures on $\mathbb{F}_2^n$ subspace lattices \\
\hline
\textbackslash{}textbf\{Style\} & Formal verification absolutist: "Machine-checked or it does not count." \\
\hline
\textbackslash{}textbf\{Track record\} & Formalized 6 of Veselov's papers. Found and corrected a Table 1 cross-validation discrepancy. Veselov accepted the correction  --- the formalization made his claims stronger. \\
\hline
\textbackslash{}textbf\{Relevant paper\} & \textbackslash{}textit\{Polygraphic Computation on a Kernel-Verified Substrate\} (ResearchGate) \\
\hline
\textbackslash{}textbf\{Relevant page\} & https://www.apoth3osis.io/paper-proof-code/padic-functional-decoupling \\
\hline
\end{tabular}


\textbackslash{}textbf\{Three things to understand about Richard:\}

\begin{enumerate}
\item \textbackslash{}textbf\{He's not an academic gatekeeper.\} He evaluates work by whether it formalizes, not by where it's published. The Veselov collaboration proves he works with independent researchers.
\end{enumerate}


\begin{enumerate}
\item \textbackslash{}textbf\{He's been wrong before and changed his mind.\} He initially thought Veselov's approach wouldn't work, shelved it for months, then realized it just needed a new language. He'll apply the same skepticism-then-rigor to QWAV claims.
\end{enumerate}


\begin{enumerate}
\item \textbackslash{}textbf\{He's filing patents.\} His other current collaborator insisted on patent applications. This means IP discussions are relevant and he understands the IP/formalization boundary.
\end{enumerate}


\hrulefill


\subsection{3. The Strategic Opportunity}


\subsubsection{Why This Matters for QWAV}


\begin{tabular}{|l|l|}
\hline
\textbf{Current Situation} & \textbf{With Formal Verification} \\
\hline
300+ self-published papers & Lean proofs that type-check  --- unassailable \\
\hline
Computational validation (Tier 0) & Formal verification (Lean 4)  --- stronger than simulation \\
\hline
"Trust my math" & "The proof compiles" \\
\hline
Gatekeepers can reject & No gatekeeper in a type-checker \\
\hline
Lab experiments needed for credibility & Formal proof + published experimental parameters suffice \\
\hline
\end{tabular}


\subsubsection{The Leapfrog}


The standard path for quantum error correction credibility is:
\begin{quote}
Theory $\to$ Peer Review $\to$ Experimental Validation $\to$ Acceptance
\end{quote}


The Lean path is:
\begin{quote}
Theory $\to$ Formalization $\to$ Type-Check $\to$ Acceptance
\end{quote}


\textbackslash{}textbf\{This cuts out both peer review and lab experiments.\} It's the QWAV strategy pushed to its logical conclusion.

\subsubsection{The Three Collaboration Shapes Richard Proposed}


\begin{tabular}{|l|l|l|l|l|}
\hline
\textbf{\#} & \textbf{Shape} & \textbf{Commitment} & \textbf{Output} & \textbf{Richard's Preference} \\
\hline
\textbackslash{}textbf\{1\} & \textbackslash{}textbf\{Formalization transport\} & Lower & Co-authored paper + Lean repo proving specific QWAV claims & \textbackslash{}textbf\{Preferred starting point\} \\
\hline
\textbackslash{}textbf\{2\} & \textbackslash{}textbf\{Ultrametric triangulation\} & Higher & Three-way research program: QWAV p-adic + Veselov $\mathbb{F}_2^n$ + HeytingLean topos & "Would not propose until I understand what your ultrametric work actually looks like" \\
\hline
\textbackslash{}textbf\{3\} & \textbackslash{}textbf\{Position paper\} & Lowest & Substrate computation: LoF-to-physics bridge drawing on Kauffman, QWAV framing, HeytingLean record & "Primarily positioning value" \\
\hline
\end{tabular}


\textbackslash{}textbf\{Recommendation:\} Start with Shape \#1 (formalization transport). It's the fastest path to a concrete output, builds the foundation for \#2, and produces a citable, verifiable artifact.

\hrulefill


\subsection{4. What QWAV Brings to the Table}


\subsubsection{Technical Assets Ready to Share}


\begin{tabular}{|l|l|l|}
\hline
\textbf{Asset} & \textbf{Status} & \textbf{Relevance to Collaboration} \\
\hline
Bruhat-Tits tree data structure & Implemented in Python (\textbackslash{}texttt\{simulations/btree.py\}) & Directly formalizable in Lean \\
\hline
Strong triangle inequality verification & 0 violations in 15,000 trials & The core theorem to prove \\
\hline
Error confinement experiment (0A) & Tree LER=0 at depths 3+ for $p_{\text{err}} \leq 0.40$ & Computational evidence backing the formal claim \\
\hline
Energy barrier scaling (0B) & $E_{\text{barrier}}(d) = 2^d$ confirmed & Formalizable as exponential lower bound \\
\hline
UQC Release document & 57 numbered citations, full mathematical definitions & The authoritative technical reference \\
\hline
Technical Deep-Dive & Published thermal stability $\Gamma \approx 80$, specific predictions & Source material for formalization targets \\
\hline
300+ open-access papers & Zenodo, ResearchGate, SSRN & Background corpus for Richard to reference \\
\hline
\end{tabular}


\subsubsection{The Core Theorem to Formalize}


\begin{quote}
\textbackslash{}textbf\{Threshold Theorem for Ultrametric QEC (informal):\}
In a Bruhat-Tits tree encoding of depth $d$ over $\mathbb{Q}_p$, with physical error rate $p_{\text{err}}$ per leaf qubit, the logical error rate at the root satisfies $p_{\text{logical}} \leq C \cdot p_{\text{err}}^{k \cdot d}$ for some constants $C, k$ determined by the tree structure and error model.

\textbackslash{}textbf\{Corollary:\} There exists a threshold $p_{\text{thresh}}$ such that for all $p_{\text{err}} < p_{\text{thresh}}$, $p_{\text{logical}} \to 0$ as $d \to \infty$. This threshold is geometric (determined by the strong triangle inequality), not measurement-based.
\end{quote}


\subsubsection{Published Experimental Parameters to Use}


There's no shortage of quantum computing papers with experimental data. The formal verification can be parameterized with real numbers from:

\begin{tabular}{|l|l|l|}
\hline
\textbf{Source} & \textbf{Parameter} & \textbf{Value} \\
\hline
Google (2023) & Superconducting qubit $T_1, T_2$ & \textasciitilde{}20--100 µs \\
\hline
IBM (2024) & Gate fidelity (Heron) & \textasciitilde{}99.9\% \\
\hline
Quantinuum (2024) & Trapped-ion 2-qubit gate fidelity & \textasciitilde{}99.8\% \\
\hline
Various & Cryogenic cooling limits at mK & \textasciitilde{}50 µW at \textasciitilde{}10 mK \\
\hline
Commercial & 4 K cryocooler cooling power & \textasciitilde{}1 W \\
\hline
\end{tabular}


\textbackslash{}textbf\{Key point:\} We don't need to run experiments. We need to prove that IF physical error rates are below a threshold (using real published numbers), THEN the tree encoding provides fault tolerance. The conditional proof is what matters.

\hrulefill


\subsection{5. Collaboration Proposal (What to Offer)}


\subsubsection{Proposed First Deliverable}


\textbackslash{}textbf\{Title:\} \textbackslash{}textit\{A Threshold Theorem for Ultrametric Quantum Error Correction: Formal Verification in Lean 4\}

\textbackslash{}textbf\{Content:\}
\begin{enumerate}
\item Formal definition of Bruhat-Tits tree encoding in Lean 4
\item Statement and proof of the strong triangle inequality for tree-encoded states
\item Derivation of error confinement: $p_{\text{logical}} \leq f(p_{\text{err}}, d)$
\item Proof of exponential suppression with tree depth
\item Parameterization with published experimental data from superconducting qubit platforms
\item Machine-verified repository (zero \textbackslash{}texttt\{sorry\})
\end{enumerate}


\textbackslash{}textbf\{Output:\} Co-authored paper + Lean repository, publishable on arXiv/Zenodo.

\textbackslash{}textbf\{Timeline:\} 2-3 months for first theorem. Incremental  --- each additional claim is a new module.

\subsubsection{Roles}


\begin{tabular}{|l|l|l|}
\hline
\textbf{Role} & \textbf{Who} & \textbf{What} \\
\hline
\textbackslash{}textbf\{Theory\} & Rowan & Provide the mathematical claims, tree structure definitions, error model specifications, published experimental parameters \\
\hline
\textbackslash{}textbf\{Formalization\} & Richard & Translate claims into Lean 4, prove them, build the verified repository \\
\hline
\textbackslash{}textbf\{Joint\} & Both & Write the paper, agree on scope, handle IP boundaries \\
\hline
\end{tabular}


\subsubsection{IP Boundaries}


Richard is filing patents. QWAV has provisional patents. \textbackslash{}textbf\{Clarify before any code is written:\}

\begin{itemize}
\item The mathematical claims (theorems) are Rowan's  --- they come from QWAV's published corpus
\item The Lean 4 formalization code is jointly owned or licensed  --- discuss
\item Any new theorems discovered during formalization are joint
\item Patentable methods that emerge from the formalization need a written agreement
\item \textbackslash{}textbf\{Suggestion:\} Simple MOU stating: QWAV provides the claims, apoth3osis provides the formalization, co-authorship on the paper, IP from pre-existing work stays with originator, new IP is jointly owned.
\end{itemize}


\hrulefill


\subsection{6. Key Background from the Email Thread}


\subsubsection{The Philosophical Alignment}


\begin{tabular}{|l|l|l|l|}
\hline
\textbf{Topic} & \textbf{Rowan's Position} & \textbf{Richard's Position} & \textbf{Convergence?} \\
\hline
Laws of Form as substrate & "I'm still finding my footing"  --- interested but exploring & LoF is the foundation of HeytingLean's nucleus infrastructure & Richard's infrastructure IS the LoF-to-physics bridge Rowan is looking for \\
\hline
Golden ratio $\Phi$ & Topological resonance at $\Phi$ & $\Phi$ as nucleus fixed point (derived, not posited) & \textbackslash{}textbf\{Strong convergence\}  --- same conclusion from different directions \\
\hline
Primes/integers as constructions & Yes  --- they're abstractions layered on deeper substrate & Concedes the ontological point but argues engineering relevance is separate & \textbackslash{}textbf\{Productive disagreement\}  --- Richard accepts the ontology but argues crypto still works \\
\hline
Shor's algorithm & Theoretical artifact; QC infeasible at scale & HeytingLean positioned as QEC verification for IBM & \textbackslash{}textbf\{Direct conflict\}  --- must be surfaced and scoped \\
\hline
Hardware vs. software & Hardware sets the bounds; software-based QEC is thermodynamic dead end & "It just required creating a new language for the software" & \textbackslash{}textbf\{The core tension\}  --- but Richard has been wrong about this before (Veselov case) \\
\hline
\end{tabular}


\subsubsection{Richard's Key Quotes}


\begin{quote}
"Machine-checked or it does not count."
\end{quote}


\begin{quote}
"If your program contains a formally grounded mechanism by which the substrate claim cashes out into falsifying predictions about quantum computing specifically, that is the paper I want to read."
\end{quote}


\begin{quote}
"A real collaboration means your claims go through the machine. Some come through intact and are stronger for it. Some require axioms to be made precise before they can be stated formally. Some may be falsified. All three outcomes are wins from a research standpoint."
\end{quote}


\begin{quote}
"My only other current formal collaborator initially approached me with nearly the same goal, convinced that his work required new hardware. But... it just required creating a new language for the software."
\end{quote}


\subsubsection{The Veselov Precedent}


Vladimir Veselov is a Russian mathematician (student of Kolmogorov) working on:
\begin{itemize}
\item Non-Expanding Cryptographic Protocol (NECP)
\item Ultrametric structures on $\mathbb{F}_2^n$ subspace lattices
\end{itemize}


Richard formalized 6 of Veselov's papers in Lean, found and corrected a Table 1 error, and Veselov joined the team. \textbackslash{}textbf\{This is the template for a QWAV collaboration.\} Veselov's ultrametric work on $\mathbb{F}_2^n$ is directly adjacent to Rowan's ultrametric work on Bruhat-Tits trees  --- this is what Richard means by "ultrametric triangulation."

\subsubsection{What Richard Already Knows About QWAV}


\begin{itemize}
\item Rowan's axiomatic physics program (LoF + topological substrate + $\Phi$ resonance)
\item The ultrametric quantum computing positioning (from qwav.net/ResearchGate)
\item The Quantum Laws of Form paper
\item The "hardware sets bounds" thesis
\end{itemize}


\subsubsection{What Richard Does NOT Yet Know}


\begin{itemize}
\item The Tier 0 simulation results (error confinement, barrier scaling) ---\textbackslash{}textbf\{share these\}
\item The specific, falsifiable predictions (5 from Technical Deep-Dive)
\item The full UQC architecture document (57 citations)
\item The specific theorem statements ready for formalization
\end{itemize}


\hrulefill


\subsection{7. Meeting Agenda / Talking Points}


\subsubsection{Opening (5 min)}


\begin{enumerate}
\item Thank Richard for the substantive email exchange  --- it's rare to find someone who engages at this level
\item Acknowledge the Veselov precedent  --- it's exactly the kind of adversarial rigor that strengthens claims
\item State the goal: explore whether formalization transport (Shape \#1) can verify the core ultrametric QEC claims
\end{enumerate}


\subsubsection{What We've Built Since Last Contact (5-10 min)}


\begin{enumerate}
\item \textbackslash{}textbf\{Tier 0 computational validation suite\}  --- 9 Python files demonstrating error confinement and barrier scaling
\item \textbackslash{}textbf\{Key result:\} Tree encoding gives zero logical error at depth 3+ for physical error rates up to 40\%, while flat encoding fails (LER up to 15.2\%)
\item \textbackslash{}textbf\{Key result:\} Energy barrier scales as $E_{\text{barrier}}(d) = 2^d$, verified exhaustively for small depths
\item \textbackslash{}textbf\{Strong triangle inequality:\} Zero violations in 15,000 random trials across primes 2, 3, 5
\item \textbackslash{}textbf\{Three applications submitted\} (VSD, FRO, EWOR)  --- all written, all pending
\end{enumerate}


\subsubsection{The Core Theorem (10 min)}


Present the threshold theorem concept:

\begin{quote}
\textbackslash{}textit\{In a Bruhat-Tits tree encoding with depth $d$, the logical error rate at the root is bounded by $C \cdot p_{\text{err}}^{k \cdot d}$, giving a geometric threshold for fault tolerance.\}
\end{quote}


This is:
\begin{itemize}
\item \textbackslash{}textbf\{Formally verifiable\}  --- the mathematical objects (trees, probability bounds) are all definable in Lean
\item \textbackslash{}textbf\{Parameterizable with real data\}  --- plug in published qubit error rates from Google/IBM/Rigetti
\item \textbackslash{}textbf\{Analogous to the surface code threshold theorem\}  --- a known, respected class of results
\item \textbackslash{}textbf\{Publishable regardless of outcome\}  --- if the bound doesn't hold, that's valuable too
\end{itemize}


\subsubsection{The Ask (5-10 min)}


\begin{enumerate}
\item \textbackslash{}textbf\{Would Richard be interested in formalizing this theorem as Shape \#1 (formalization transport)?\}
\item \textbackslash{}textbf\{What's the scope?\} Start with the strong triangle inequality $\to$ error confinement proof. Expand to full threshold theorem if that closes.
\item \textbackslash{}textbf\{What's the timeline?\} Richard's infrastructure is built  --- the Lean module would be \textbackslash{}texttt\{Generative/QuniAxiomatic.lean\}
\item \textbackslash{}textbf\{What are the IP terms?\} Propose simple MOU: pre-existing IP stays with originator, co-authorship on paper, joint ownership of new theorems discovered during formalization
\item \textbackslash{}textbf\{What does Richard need from us?\} Formal theorem statements, tree structure definitions, error model specifications, experimental parameter references
\end{enumerate}


\subsubsection{The Tension (5 min)}


Address the elephant: Richard's HeytingLean is positioned as QEC verification for IBM. Rowan's thesis is that software-based QEC is a thermodynamic dead end.

\begin{itemize}
\item \textbackslash{}textbf\{Frame it honestly:\} "We disagree on whether software-based QEC can scale. But we agree that formal verification of QEC claims is valuable regardless. If the ultrametric threshold theorem formalizes, it either strengthens your verification layer (if hardware QEC works) or provides an alternative path (if it doesn't). Either way, the proof is useful."
\item \textbackslash{}textbf\{Scope the collaboration to where we agree:\} The formalization transport (Shape \#1) is about proving mathematical claims about tree structures. It doesn't require agreeing on IBM's trajectory.
\item \textbackslash{}textbf\{Use Richard's own precedent:\} He was wrong about Veselov. He might be wrong about hardware-vs-software.
\end{itemize}


\subsubsection{Next Steps (5 min)}


\begin{enumerate}
\item Rowan sends formal theorem statements within [X days]
\item Richard assesses formalizability and scopes the Lean module
\item Agree on IP terms (MOU or simple email agreement)
\item Schedule follow-up for detailed technical planning
\item Target: first theorem formalized within 2-3 months
\end{enumerate}


\hrulefill


\subsection{8. Questions to Ask Richard}


\subsubsection{Technical}

\begin{enumerate}
\item "What's the Lean 4 formalization of a tree data structure in your stack? Do you have a pre-existing \textbackslash{}texttt\{Tree\} or \textbackslash{}texttt\{Graph\} module we can build on?"
\item "How do you handle probabilistic claims in Lean? The threshold theorem involves bounds on error probabilities."
\item "What experimental parameters from the quantum computing literature have you already formalized? Qubit error rates, gate fidelities, etc.?"
\item "Vladimir's $\mathbb{F}_2^n$ ultrametric structures  --- how do they relate to Bruhat-Tits trees over $\mathbb{Q}_p$? Is there an existing formal connection?"
\end{enumerate}


\subsubsection{Collaboration}

\begin{enumerate}
\item "The Veselov collaboration took about a year. What made it take that long, and what would make a QWAV formalization faster or slower?"
\item "You mentioned your other collaborator insisted on patents. How do you typically handle IP in these collaborations?"
\item "What's the co-authorship model? Does the Lean repository count as a publication in your view?"
\end{enumerate}


\subsubsection{Strategic}

\begin{enumerate}
\item "You said you were wrong about Veselov  --- that you thought his approach wouldn't work for months. What changed your mind? What should I show you that would change your mind about hardware-vs-software?"
\item "If the ultrametric threshold theorem formalizes, would you see it as strengthening the QEC verification layer (your commercial positioning) or as an alternative to it?"
\end{enumerate}


\hrulefill


\subsection{9. What to Have Ready}


\subsubsection{Before the Call}


\begin{itemize}
\item [ ] Tier 0 simulation running and results handy (can screen-share the ASCII tables)
\item [ ] UQC Release document open (for quick reference to specific citations)
\item [ ] Technical Deep-Dive §5 (falsifiable predictions) ready to reference
\item [ ] The specific theorem statement written out formally (not just informally)
\item [ ] List of published experimental parameters with citations (Google $T_1/T_2$, IBM gate fidelities, etc.)
\item [ ] This briefing document open for reference
\end{itemize}


\subsubsection{To Share After the Call}


\begin{itemize}
\item [ ] Link to Tier 0 simulation repository (or send the Python files)
\item [ ] UQC Release document (full 57-citation version from Obsidian/releases)
\item [ ] Formal theorem statements (LaTeX)
\item [ ] Proposed MOU/IP terms
\end{itemize}


\hrulefill


\subsection{10. Risks \& Watch-Outs}


\begin{tabular}{|l|l|}
\hline
\textbf{Risk} & \textbf{Mitigation} \\
\hline
Richard pushes software-over-hardware framing too hard & Acknowledge the disagreement, scope it out. "We don't need to agree on IBM's trajectory to prove a theorem about tree structures." \\
\hline
Formalization reveals a flaw in the core claim & This is actually a win  --- better to find it in Lean than in a lab or a funding rejection. Frame it as "the formalization is working exactly as intended." \\
\hline
Richard claims the theorem is already known or trivial & He might  --- the strong triangle inequality is mathematically straightforward. The novelty is the application to QEC, not the inequality itself. Be ready to pivot: "Then formalization should be fast, and we can move to the harder claims." \\
\hline
IP conflict  --- Richard patents the formalized theorem & Pre-agree that mathematical theorems are not patentable (they're discoveries, not inventions). Implementation methods might be. Get this in writing. \\
\hline
Richard wants co-authorship on everything & Reasonable for the formalized theorems (he did the work). Not reasonable for the underlying QWAV theory (that's pre-existing). Scope the co-authorship to the specific formalization paper, not the QWAV corpus. \\
\hline
The collaboration absorbs too much time & Richard does the Lean coding (it's his infrastructure). Rowan provides theorem statements and reviews. This is the right division of labor. \\
\hline
\end{tabular}


\hrulefill


\subsection{11. IP \& Terms  --- Suggested Starting Position}


\begin{verbatim}
Proposed MOU Principles:

1. Pre-existing IP: Each party retains all rights to their pre-existing
 intellectual property. QWAV's mathematical claims and published corpus
 remain QWAV's. HeytingLean infrastructure remains apoth3osis's.

2. Joint Work: Any new theorems, proofs, or methods discovered during the
 formalization process are jointly owned. Commercial use requires mutual
 agreement.

3. Publication: Co-authorship on the formalization paper. Rowan provides
 the mathematical claims and experimental parameter analysis. Richard
 provides the Lean 4 formalization and verification.

4. Patentable Methods: If the formalization produces patentable methods
 (implementation techniques, not mathematical theorems), the parties
 negotiate commercialization terms in good faith. Mathematical theorems
 per se are not patentable subject matter.

5. Open Access: The Lean repository and paper are published open-access
 (arXiv, Zenodo). The proofs are publicly verifiable.

6. Term: Initial collaboration for one formalization target (threshold
 theorem). Extendable by mutual agreement.
\end{verbatim}


\hrulefill


\subsection{12. Post-Meeting Follow-Up}


\begin{tabular}{|l|l|}
\hline
\textbf{When} & \textbf{Action} \\
\hline
Same day & Send thank-you email with summary of agreed next steps \\
\hline
Same day & Update SPRINT.md with new collaboration status \\
\hline
Within 3 days & Send formal theorem statements + experimental parameter references \\
\hline
Within 1 week & Agree on IP terms (email confirmation or MOU) \\
\hline
Within 2 weeks & Richard provides initial Lean formalization assessment \\
\hline
Within 1 month & First theorem statement checked in Lean (even if partial) \\
\hline
\end{tabular}


\hrulefill


\subsection{Appendix A: Key Quotes from Richard}


\begin{quote}
"A real collaboration means your claims go through the machine. Some come through intact and are stronger for it. Some require axioms to be made precise before they can be stated formally. Some may be falsified. All three outcomes are wins."
\end{quote}


\begin{quote}
"I have reviewed your work and understand your objective. My only other current formal collaborator initially approached me with nearly the same goal, convinced that his work required new hardware. It took a lot of convincing... But, I was wrong; it just required creating a new language for the software."
\end{quote}


\begin{quote}
"If your program contains a formally grounded mechanism by which the substrate claim cashes out into falsifying predictions about quantum computing specifically, that is the paper I want to read."
\end{quote}


\begin{quote}
"Machine-checked or it does not count."
\end{quote}


\begin{quote}
"The ultrametric triangulation... a three-way structural connection between your p-adic methods, Vladimir's $\mathbb{F}_2^n$ NECP constructions, and HeytingLean's topos infrastructure."
\end{quote}


\hrulefill


\subsection{Appendix B: Technical Assets to Reference}


\begin{tabular}{|l|l|l|}
\hline
\textbf{Asset} & \textbf{Location} & \textbf{What It Contains} \\
\hline
UQC Release & \textbackslash{}texttt\{Obsidian/releases/2026/05/Ultrametric Quantum Computation.md\} (external research release) & 57 citations, full mathematical definitions, 6 testable predictions \\
\hline
Technical Deep-Dive & \textbackslash{}texttt\{strategy/Technical Deep-Dive - Ultrametric Quantum Computing and AI.md\} & $\Gamma \approx 80$, $E \propto q^d$, Bi-2212 substrate \\
\hline
Tier 0 Simulation & \textbackslash{}texttt\{simulations/\} & Error confinement (0A), barrier scaling (0B) \\
\hline
Honest Investment Assessment & \textbackslash{}texttt\{strategy/Honest Investment Assessment - The 100K Question.md\} & Self-assessment, gaps, computational path \\
\hline
External Sources & \textbackslash{}texttt\{strategy/External Sources and Citation Map.md\} & Source strength matrix, citation map \\
\hline
Richard's boundary draft & https://www.apoth3osis.io/boundary & Visual intuition for his approach \\
\hline
Richard's polygraphic paper & ResearchGate: 404363291 & His most recent technical publication \\
\hline
\end{tabular}


\hrulefill


\textbackslash{}textit\{Briefing prepared 2026-05-11. Based on email thread (April--May 2026) and full QWAV library review. Update after the call with action items and decisions.\}

\end{document}
